\documentclass[12pt,a4paper,oneside]{article}
\usepackage[brazil]{babel}
\usepackage[utf8]{inputenc}
\usepackage{olymp}
\usepackage{color}
\usepackage[dvipsnames]{xcolor}
\usepackage{listings}

\setcounter{problem}{4}

\begin{document}

\begin{problem}{Números espelhos}{espelho}{2}

Dois números inteiros $A$ e $B$ são ditos espelhos se ao inverter-se a disposição dos algarismos de $A$ obtêm-se o valor de $B$. Por exemplo, 5273 e 3725 são espelhos.

Elabore um programa para verificar se dois números inteiros de 4 dígitos são espelhos entre si. Para resolução, seu programa deve usar obrigatoriamente uma função com o seguinte
protótipo:

\texttt{\textcolor{Mulberry}{int} espelhar(\textcolor{Mulberry}{int} n)}

A função recebe como parâmetro um número inteiro de 4 dígitos e retorna-o espelhado.

%\Notes
%
%Observações, se tiver.

\InputFile

A entrada contém vários casos de teste. Cada caso de teste é dado em uma linha contendo dois valores inteiros $A$ e $B$, separados por um espaço em branco. Restrições: $1000 \leq A, B \leq 9999$.

A entrada termina com $A = 0$ e $B = 0$, que não deve ser processada.

\OutputFile

A saída consiste em várias linhas, uma para cada caso de teste da entrada. Na linha deverá aparecer a palavra ``\texttt{espelho}'', caso os valores entrados sejam espelho entre si, ou ``\texttt{nao espelho}'', caso contrário.

% \Notes
% Nesta questão, seu programa deve usar a função \texttt{main} dada acima exatamente como está, não a modifique! Sua tarefa é apenas criar a função \texttt{eh\_primo}.

\Example

\begin{example}
\exmp{
6946 4532
5786 6875
7544 3754
5427 7245
9746 6479
0 0
}{
nao espelho
espelho
nao espelho
espelho
espelho
}%
\end{example}



%Comentários sobre o exemplo, se necessário.

\end{problem}

\end{document}
